\chapter{Groin Pain or Swelling}

\section*{Definition}
Pain, discomfort, or visible enlargement in the inguinal or pubic region, arising from musculoskeletal, hernia-related, vascular, or genitourinary causes.

\section*{Pathophysiology}
Groin pain in younger patients commonly results from adductor muscle strain (sports hernia or athletic pubalgia). Hernias present with bulging aggravated by increased intra-abdominal pressure. Lymphadenopathy suggests infection or malignancy. Testicular causes (torsion, epididymitis) refer pain to groin via shared innervation.

\section*{Etiology}
\begin{itemize}
  \item Inguinal or femoral hernia
  \item Adductor muscle strain (groin pull)
  \item Lymphadenopathy (inguinal lymph nodes)
  \item Epididymitis or orchitis
  \item Testicular torsion or tumor
  \item Osteitis pubis
  \item Hip joint pathology (labral tear, OA)
  \item Kidney stones (referred pain)
  \item Femoral artery pseudoaneurysm (post-catheterization)
  \item Psoas abscess
\end{itemize}

\section*{Indications for Evaluation}
Seek care if:
\begin{itemize}
  \item Severe or worsening symptoms
  \item Severe groin pain with nausea/vomiting (possible strangulated hernia), testicular pain with swelling, or fever with groin pain
  \item Accompanied by other concerning signs
\end{itemize}

\section*{Related Symptoms}
\begin{itemize}
  \item Hip pain
  \item Lower abdominal pain
  \item Testicular pain
  \item Nausea
  \item Lump in groin
\end{itemize}
\textbf{Note:} Always consider serious underlying conditions when evaluating persistent groin pain or swelling.

\section*{Self-Care / Home Management}
\begin{itemize}
  \item Rest and avoid activities that may aggravate the condition.
  \item Maintain adequate hydration and a balanced diet.
  \item Over-the-counter remedies may provide symptomatic relief (consult a pharmacist or doctor).
  \item Monitor symptoms and seek professional advice if they persist or worsen.
  \item Avoid self-medicating with prescription medications without medical guidance.
\end{itemize}

\section*{Diagnostic Approach}
\begin{itemize}
  \item A thorough history and physical examination are the first steps in evaluation.
  \item Laboratory tests (blood work, urinalysis) may be ordered based on clinical suspicion.
  \item Imaging studies (X-ray, ultrasound, CT, MRI) may be indicated when structural pathology is suspected.
  \item Specialist referral is arranged when the diagnosis remains uncertain or requires specific expertise.
\end{itemize}

\section*{Treatment / Management Overview}
\begin{itemize}
  \item Treatment targets the underlying cause identified through diagnostic evaluation.
  \item Supportive care and symptom management are provided while awaiting definitive therapy.
  \item Medications, lifestyle modifications, and physical therapies may be prescribed as appropriate.
  \item Follow-up ensures treatment effectiveness and allows timely adjustments to the care plan.
\end{itemize}

\clearpage